\section{Annexes}

\begin{guideline}
Les annexes constituent la base probatoire de votre dossier. Chaque trace doit être accompagnée d'un commentaire qui explique son contexte et ce qu'elle démontre. Organisez-les par compétence pour faciliter la lecture des évaluateurs. Au moins 3 traces de natures différentes par compétence sont attendues.
\end{guideline}


\subsection{Liste des traces par compétence}

\begin{guideline}
Pour chaque compétence visée, listez les traces associées (documents, captures, extraits de code, tableaux...) avec leur référence (ex. : Annexe A.1), une description courte, et un commentaire expliquant en quoi cette trace démontre l'acquisition de la compétence. Pensez à varier les types de traces.
\end{guideline}

\begin{center}
\renewcommand{\arraystretch}{1.5}
\begin{tabularx}{\textwidth}{@{}l l l X@{}}
\toprule
\textbf{Réf.} & \textbf{Type de trace} & \textbf{Description} & \textbf{Commentaire --- compétence(s) et composantes démontrées} \\
\midrule
\multicolumn{4}{l}{\textbf{Compétence 1 : \instruction{[Nom de la compétence]}}} \\
A.1 & \instruction{Document} & \instruction{Cahier des charges rédigé} & \instruction{Expliquez ce que cette trace démontre, quelles composantes sont adressées...} \\
A.2 & \instruction{Capture} & \instruction{Résultat automatisé} & \instruction{Commentaire associé...} \\
A.3 & \instruction{Livrable} & \instruction{Extrait de code commenté} & \instruction{Commentaire associé...} \\
\bottomrule
\end{tabularx}
\end{center}

\begin{guideline}
Rappel : au moins 3 traces de natures différentes sont nécessaires par compétence pour en valider l'acquisition. Reproduisez ce tableau autant de fois que nécessaire, une fois par compétence visée.
\end{guideline}


\subsection{Évaluation du maître de stage}

\begin{guideline}
Ce document est obligatoire. Il doit être préparé par vous (stagiaire) sous forme d'une synthèse de votre échange de mi-stage ou de fin de stage avec votre maître de stage, puis contresigné par lui/elle. Il atteste de votre niveau d'avancement et des retours reçus.
\end{guideline}

\instruction{[Insérez ici la synthèse de l'échange avec votre maître de stage, préparée par vos soins. Elle doit couvrir : l'avancement de la mission, les points forts identifiés, les axes d'amélioration, et les retours sur votre posture professionnelle. Ce document doit être contresigné par votre maître de stage.]}


% ==========================================================
% Rappel des principes directeurs
% ==========================================================
\newpage
\subsection*{Rappel : principes directeurs}

\begin{guideline}
Ces principes s'appliquent à l'ensemble de votre dossier. Relisez-les avant de finaliser chaque section.
\end{guideline}

\begin{itemize}
    \item \textbf{Écrire pour démontrer} : Chaque partie doit établir un lien explicite entre vos activités et les compétences visées.
    \item \textbf{Contribution personnelle} : Distinguez clairement vos actions propres de celles de l'équipe. Évitez le « nous » sans précision.
    \item \textbf{Preuves systématiques} : Chaque affirmation doit s'appuyer sur une trace commentée en annexe.
    \item \textbf{Posture réflexive} : Analysez vos choix, assumez vos limites, proposez des alternatives ou des recommandations.
    \item \textbf{Qualité rédactionnelle} : Soignez la clarté, les références aux annexes et la cohérence visuelle du document.
\end{itemize}

\vspace{1cm}
\begin{center}
\instruction{--- Les éléments en magenta italic sont des consignes à supprimer avant de rendre le dossier ---}
\end{center}
